\scanpage{30}
\mspage{29}
SYSTÈMES de RACINES (10) \hfill Tribu n° 59

un quotient sur un anneau $A$ qqconque ; ce quotient est unique car $d$ n'est pas diviseur de zéro. Pour (ii), on écrit $d=\sum_x c_x Qe^x$, (cf. prop. 7) ; or, dans $d$, il y a un terme en $e^{w_0}$, et tous les autres sont plus petits, (pour l'ordre défini par $P^+$) ; comme les $x$ sont $\geq w_0$, il y a un seul terme dans la somme, avec coeff. 1.

Le haut de la p. 162, (3 lignes), aura déjà été dit.

§ 8. \uuline{Transformations de Coq-se-taire.}

\uline{n° 1} - Le lemme 6 du § 6, n° 13 remonte ici.

\uline{n° 2} - Dans le 2) en bas de la p. 163, prendre seulement $g=s_1$. Dans le lemme 1, on notera $L$ le monoïde, $Y$ l'ensemble des $n!$ mots et $C(X)$ le graphe. Dans l'énoncé, remplacer 2) par : `` y'a un mot $A$ et un sommet $b$ tels que $u=Ab$ et $v=bA$ '' ; ça simplifie la démonstration. Mieux expliquer comment la connexion du graphe donne la conjugaison voulue.

La fin de la démonstration, (p. 165), se simplifie ainsi : il y a deux cas ; si $a$ et $b$ sont non liés, c'est facile comme dans le texte ; si $a$ et $b$ sont liés, alors $a$ n'est lié à aucun élément de $A$, et on passe de $bAa$ à $Aba$ par les étapes $bAa$, $baA$, $aAb$, $Aba$.

\uline{n° 3} - Commencer par un laïus introduisant les $m_i$, (avec déf. en forme des exposants), et donnant les formules (1), (2), (3), (numérotées), qui sont des trivialités. Puis une proposition, réduite à (ii) et (iii), où on supposera $G$ irréductible, ((ii) est faux avec $G$ réductible). Pour la démonstration, on
